
\medskip
\phantomsection\label{rem:sga3-VIB-6.3}%
\noindent\textbf{Remark 6.3.---}
\emph{Let \(k\) be an algebraically closed field, let \(G\) be a
\(k\)-group, and let \(H\) be a group subscheme of \(G\). Suppose
that \(G\) and \(H\) are of finite type over \(k\) and are
reduced. Then
\(\underline{\operatorname{Norm}}_G H\) (respectively,
\(\underline{\operatorname{Centr}}_G H\)) is representable by a
group subscheme of \(G\), whose associated reduced subscheme is
precisely the normalizer (respectively, centralizer) of \(H\) in
\(G\) in the sense of the \emph{Bible}.}

\medskip
\phantomsection\label{prop:sga3-VIB-6.4}%
\noindent\textbf{Proposition 6.4.---}
\emph{Let \(G\) be an \(S\)-group, and let \(u:X\to G\) be a
monomorphism of \(S\)-schemes. Set}
\[
  T=\underline{\operatorname{Transp}}_G(X,X).
\]
\emph{The following conditions are equivalent:}
\begin{enumerate}
\item[\textup{(i)}]\phantomsection\label{cond:sga3-VIB-6.4-i}
  \emph{\(T\) is a group sub-\(S\)-functor of \(G\).}
\item[\textup{(ii)}]\phantomsection\label{cond:sga3-VIB-6.4-ii}
  \emph{\[
    T=\underline{\operatorname{Transp}}^{\mathrm{str}}_G(X,X)
     =\underline{\operatorname{Norm}}_G X.
  \]}
\end{enumerate}
\emph{These conditions hold in each of the following two cases:}
\begin{enumerate}
\item[\textup{a)}]\phantomsection\label{case:sga3-VIB-6.4-a}
  \emph{\(X\) is of finite presentation over \(S\).}
\item[\textup{b)}]\phantomsection\label{case:sga3-VIB-6.4-b}
  \emph{\(T\) is representable by a scheme of finite
  presentation over \(S\).}
\end{enumerate}

The equivalence of \textup{(i)} and \textup{(ii)} follows from
the fact that, for every morphism \(S'\to S\) and all
\(t,t'\in T(S')\), one has
\[
  tt'\in T(S'),
\]
and from the identity
\[
  \underline{\operatorname{Transp}}^{\mathrm{str}}_G(X,X)
  =T\cap c(T)
\]
(cf.\ \hyperref[def:sga3-VIB-6.1]{6.1\textup{(ii)}}).

Suppose first that \textup{a)} holds. If \(t\in T(S)\), then
\(\operatorname{int}(t)\) is a monomorphism from \(X\) to itself,
and therefore an \(S\)-automorphism of \(X\) by EGA~IV\(_4\),
17.9.6. Thus
\[
  t\in\underline{\operatorname{Transp}}^{\mathrm{str}}_G(X,X),
\]
which proves the assertion in case \textup{a)}.

In case \textup{b)}, it is clear that
\[
  \mu(T\times_S T)\subset T,
\]
and the assertion follows from the next lemma.

\medskip
\phantomsection\label{lem:sga3-VIB-6.4.2}%
\noindent\textbf{Lemma 6.4.2.---}%
\emph{Let \(G\) be an \(S\)-scheme of finite presentation,
equipped with an associative law (in the sense of
EGA~0\(_{\mathrm{III}}\), 8.2.5). Suppose that, for every
\(S\)-scheme \(S'\) and every \(g\in G(S')\), right and left
translation by \(g\) in the set \(G(S')\) are injective, and
suppose that \(G(S)\ne\varnothing\). Then \(G\) is an
\(S\)-group.}

\emph{Proof.}
It is enough to prove that \(G(S')\) is a group for every
\(S\)-scheme \(S'\). The hypothesis immediately implies that
right and left translation by every \(g\in G(S')\) on
\(G_{S'}\) are \(S'\)-monomorphisms
\[
  G_{S'}\longrightarrow G_{S'}.
\]
They are therefore \(S'\)-automorphisms, since \(G\) is of
finite presentation over \(S\) (EGA~IV\(_4\), 17.9.6).
Consequently, right and left translation by \(g\) in the set
\(G(S')\) are bijective, and the assertion reduces to the
following lemma.

\medskip
\phantomsection\label{lem:sga3-VIB-6.4.3}%
\noindent\textbf{Lemma 6.4.3.---}
\emph{Let \(G\) be a nonempty set with an associative law such
that right and left translations are bijective. Then \(G\) is a
group.}

The proof is left to the reader.

\medskip
\phantomsection\label{def:sga3-VIB-6.5}%
\noindent\textbf{Definition 6.5.---}
\emph{Let \(G\) be an \(S\)-group, let \(H\) be an
\(S\)-functor, and let \(u:H\to G\) be a monomorphism.}
\begin{enumerate}
\item[\textup{(i)}]
  The \emph{connected centralizer of \(H\) in \(G\)}, denoted
  \[
    \underline{\operatorname{Centr}}^0_G H,
  \]
  is the identity component of the functor
  \(\underline{\operatorname{Centr}}_G H\)
  (cf.\ \hyperref[def:sga3-VIB-3.1]{3.1} and
  \hyperref[def:sga3-VIB-6.1]{6.1\textup{(iii)}}).
\item[\textup{(i\(^{\prime}\))}]
  For every morphism \(u:X\to G\), one similarly defines the
  functor
  \[
    \underline{\operatorname{Centr}}^0(u)
  \]
  (cf.\ \SGARefLink{prop:sga3-VIB-6.2}{6.2}\textup{(iv)}).
\item[\textup{(ii)}]
  When \(u_s\) is a closed immersion for every \(s\in S\), the
  \emph{connected normalizer of \(H\) in \(G\)}, denoted
  \[
    \underline{\operatorname{Norm}}^0_G H,
  \]
  is the identity component of the functor
  \(\underline{\operatorname{Norm}}_G H\).
\end{enumerate}

\noindent\emph{N.B.}
By \hyperref[cor:sga3-VIB-1.4.2]{1.4.2}, the hypothesis in
\textup{\SGARefLink{cor:sga3-VIB-1.4.2:item:ii}{(ii)}} is satisfied when \(H\) is an \(S\)-group scheme,
\(G\) and \(H\) are locally of finite type over \(S\), and
\(u\) is a quasi-compact morphism of \(S\)-groups.
